\begin{proof}

\begin{enumerate}[label={$\langle1\rangle$\arabic*.},leftmargin=*]
\item We resolve the pair by a proper holomorphic modification
$\pi:U'\to U$ which is an isomorphism away from $S$ and over
$S_{\mathrm{reg}}$. The total central divisor has normal crossings.
The strict transform has multiplicity $m$, and each exceptional
component has multiplicity at least $2m$. Indeed, its first centre
lies over a singular point of the reduced divisor, whose order is
at least two; subsequent centres either retain such an exceptional
multiplicity or add further positive multiplicities.

\item We spell out the trace comparison. At a point on $k$ branches,
with local equation $z_1^{n_1}\cdots z_k^{n_k}=t$, the nearby
fibre retracts onto a disjoint union of real $(k-1)$-tori.
Monodromy permutes these components and translates any component
that it preserves. For $k\ge2$ the alternating cohomological
trace is zero. On a single branch of multiplicity $n$ there are
$n$ contractible nearby components in a cyclic permutation;
the $d$th iterate has trace $n$ when $n\mid d$, and zero otherwise.
Alternating traces in the nearby-cycle spectral sequence, with a
finite stratification of the compact central divisor, therefore give
\[
 L(A^d;H^*(T^4;\C))=\sum_{n_j\mid d}n_j e(E_j^\circ).
\]
This is the global form of A'Campo 1975, Theorem 1; his proper
analytic construction in Section 2 and spectral-sequence proof in
Section 3 give precisely this comparison. The explicit local
calculation verifies the hypotheses without a Kähler assumption.
For $d=m$ the strict transform is the sole contribution, and its
open stratum is $S_{\mathrm{reg}}$.

\item These same charts prove quasi-unipotence. A common multiple $N$
of the $n_j$ makes monodromy act trivially on all local
nearby-cohomology sheaves, hence on the $E_2$ page. Its action
on the filtered abutment is consequently unipotent. Every
eigenvalue of $A$ is a root of unity. Since torus cohomology is
the exterior algebra on $H_1$, its Lefschetz number is
$\det(I-A^m)$. Complex conjugate eigenvalues contribute
$|1-\lambda^m|^2$; real eigenvalues are $1$ or $-1$ and
contribute zero or two. Thus this determinant is nonnegative
and at most $2^4$, as required.

\end{enumerate}
\end{proof}
